% Reproduced inherited proof; finding P-F28.
\subsection{The exact local coefficient of the chain component}

\begin{proposition}[P-F28: chain multiplicity and cyclotomic valuation]
\label{prop:chainlocalcoefficient}
Let $n\ge4$, $M=2^{n-1}$, $g=\gcd(n,3)$, and
$m=\binom{n+2}{3}$. Let $Y\subset X=\operatorname{Spec}T^{(n)}$
be the reduced chain orbit closure of P-F24. Its grading
stabilizer at every nonzero point is exactly
$\boldsymbol\mu_M$. In the notation of P-F23 its leading
graded multiplicity is
\[
 e_Y=\lim_{z\to1}(1-z)
           \sum_{s\ge0}\dim\mathbb C[Y]_{Ms}z^s=1.
 \tag{P-F28a}
\]

Fix a primitive $M$th root $\zeta$, and let $g_\zeta$
denote its grading action. Normalize the cyclotomic
valuation above $2$ by $v(1-\zeta)=1$. The generic
conormal representation of $Y$ in $X$ has dimension
$m-n^2$, and
\[
 v\!\left(\det(1-g_\zeta\mid N^*_{Y/X,\eta_Y})\right)
       =3M-\binom{n+2}{2}.
 \tag{P-F28b}
\]
Consequently the local coefficient attached to this
component by the calculation of P-F23,
\[
 \mathcal C_n(\zeta):=
   \frac{e_Y}{\det(1-g_\zeta\mid N^*_{Y/X,\eta_Y})},
\]
is nonzero and satisfies
\[
 v(\mathcal C_n(\zeta))=\binom{n+2}{2}-3M.
 \tag{P-F28c}
\]
This is a local coefficient. No claim is made that it
equals a Laurent coefficient of the entire Hilbert series,
or dominates the contributions of other possible fixed components
by this proposition alone. P-F30--P-F31 below prove that it is
the actual normalized leading coefficient.
\end{proposition}

\begin{proof}
Put $\gamma=3/g$. By P-F24, the full group of linear
transformations fixing the chain form exactly is cyclic
of order $3M$, and its determinant image has order
$3M/g$. On the quotient, a matrix $B\in GL_n$ acts by
the grading parameter $(\det B)^{-\gamma}$, with the
point-action convention of (P-F26b). Thus the image
of this exact-form stabilizer in the grading torus has
order
\[
 \frac{3M/g}{\gcd(3M/g,\,3/g)}=M.
 \tag{P-F28d}
\]
This image is the entire grading stabilizer of the
quotient point. To see the converse inclusion, let a
grading parameter $u$ fix that point and choose a
matrix $B$ such that $(\det B)^{-\gamma}=u$.
The chain form and $B$ applied to it have closed
$SL_n$ orbits and the same quotient point. The
unique-closed-orbit property of the affine reductive
quotient therefore gives $S\in SL_n$ such that
$B\cdot F_n=S\cdot F_n$. The matrix $S^{-1}B$
fixes the form exactly and induces the same parameter
$u$. This proves the converse and hence the exact
stabilizer assertion. Grading translations give it
at every point of the same nonzero orbit.

Restrict homogeneous functions to the map
$\mathbb G_m\to Y$, $u\mapsto u\cdot p_n$.
A function of degree $d$ restricts to $P(p_n)u^d$.
The image of the coordinate ring is consequently
\[
 \mathbb C[Y]=\mathbb C[u^d:d\in S],
 \tag{P-F28e}
\]
where $S\subset\mathbb Z_{\ge0}$ is a finitely generated
additive semigroup containing a positive integer.
Finite generation follows by restricting a finite
set of homogeneous algebra generators; positivity
follows, for example, from the nonzero discriminant.
The common stabilizer of these monomials is
$\boldsymbol\mu_{\gcd S}$, so (P-F28d) gives
$\gcd S=M$.
This description also shows that there are no other
nonzero points of $Y$: the map $\mathbb A^1\to Y$
from (P-F28e) is finite and surjective, since $u$
is integral over the ring containing some $u^a$,
and its points with $u\ne0$ are exactly the grading
orbit.

After dividing its degrees by $M$, this semigroup
has greatest common divisor one and contains every
sufficiently large integer. Here is the elementary
argument for the latter assertion. Choose a positive
element $a$ of the semigroup. The residues of the
semigroup elements generate all of the finite group
$\mathbb Z/a\mathbb Z$, because their greatest common
divisor with $a$ is one. A subsemigroup of a finite
group is a subgroup: the inverse of an element of
order $b$ is its $(b-1)$-fold sum. Hence there is a
semigroup element in every residue class modulo $a$.
Adding nonnegative multiples of $a$ to these finitely
many representatives proves the assertion. Each
nonzero graded piece in (P-F28e) is one-dimensional.
The regraded Hilbert series therefore differs from
$1/(1-z)$ by a polynomial. This proves $e_Y=1$.

We next compute the conormal valuation. Write
$r=n-1$, so $M=2^r$. For any nonzero residue
$b\bmod M$, one has
\[
 v(1-\zeta^b)=2^{v_2(b)},
 \qquad 0\le v_2(b)\le r-1.
 \tag{P-F28f}
\]
For completeness, the polynomial
$\Phi_{2^r}(1+Z)=(1+Z)^{2^{r-1}}+1$
is Eisenstein at $2$: its constant term is $2$,
and reduction modulo $2$ leaves only the leading
term. Thus $1-\zeta$ is a uniformizer of the
cyclotomic extension at $2$.
If $b=2^j c$ with $c$ odd, the quotient
$(1-\zeta^{2^j c})/(1-\zeta^{2^j})$ is a unit,
since its geometric sum is congruent to the odd
integer $c$ modulo the uniformizer; one may choose
$c$ positive in its residue class. Moreover
\[
 1-\zeta^{2^j}
   =\prod_{k=0}^{2^j-1}
       \left(1-\zeta^{\,1+(M/2^j)k}\right).
 \tag{P-F28g}
\]
Every exponent on the right is odd, and every factor
has valuation one. This gives (P-F28f).
In particular this valuation is unchanged on
replacing $\zeta$ by any primitive $M$th root,
or on replacing a character by its inverse.

Choose a primitive $M$th root $\xi$ so that the
diagonal chain automorphism
\[
 A=\operatorname{diag}
    (\xi,\xi^{-2},\xi^4,\ldots,\xi^{(-2)^{n-1}})
\]
induces $g_\zeta$ on the quotient. This is possible:
its determinant exponent
$S_0=(1-(-2)^n)/3$ and the integer $\gamma$ are
both odd, so the quotient parameter
$\xi^{-\gamma S_0}$ runs through all primitive
$M$th roots. Put $w_i=(-2)^i\bmod M$ for
$0\le i\le n-1$. The weights $w_0,\ldots,w_{n-2}$
are nonzero and have respective two-adic valuations
$0,\ldots,n-2$, while $w_{n-1}=0$.

By P-F24, the tangent space of $X$ at the chain point
is the quotient of the cubic representation by the
tangent space of the $SL_n$ orbit. Its fixed part is
one-dimensional, namely the tangent direction of
the grading orbit. Taking moving parts and then
duals therefore gives the conormal representation
of $Y$. The cubic space has dimension $m$, the
$SL_n$ orbit has dimension $n^2-1$, and the fixed
quotient direction has dimension one. The conormal
dimension is consequently $m-n^2$.
The same character data holds at the generic point
of $Y$, since its nonzero grading orbit is dense
and the grading torus commutes with $g_\zeta$.

It suffices to subtract the determinant valuation
of the moving adjoint representation from that
of the moving cubic representation. Duals do not
change these valuations by (P-F28f).
The off-diagonal adjoint weights are $w_i-w_j$
for $i\ne j$. If $i<j$, their two-adic valuation
is $i$, including the case $j=n-1$ of the zero
weight. The adjoint determinant valuation is thus
\[
 V_{\mathrm{ad}}
   =2\sum_{i=0}^{n-2}(n-1-i)2^i
   =4M-2n-2.
 \tag{P-F28h}
\]

For cubic monomials, fix the smallest index $i$
of a nonzero weight that occurs. There are three
cases according to its multiplicity. If it occurs
once, the other two indices are greater than $i$,
and there are $\binom{n-i}{2}$ unordered choices
with repetition. The weight sum has two-adic
valuation $i$. If it occurs three times, there is
one monomial, also with valuation $i$.
If it occurs twice, the remaining index is some
$j>i$. The choice $j=i+1$ is a fixed chain
monomial and supplies no moving factor. Each
$j\ge i+2$ instead gives weight $2w_i+w_j$ of
valuation $i+1$; there are $n-i-2$ such choices
for $0\le i\le n-3$. The final all-zero-weight
cube is fixed. These cases exhaust the cubic
monomials, so
\begin{align*}
 V_{\mathrm{cub}}
 &=\sum_{i=0}^{n-2}
        \left[\binom{n-i}{2}+1\right]2^i
      +\sum_{i=0}^{n-3}(n-i-2)2^{i+1}\\
 &=7M-\frac{n^2+7n+6}{2}.
 \tag{P-F28i}
\end{align*}
For the finite sums used here, the identities are
\[
 \sum_{i=0}^{n-2}(n-1-i)2^i=2^n-n-1,\qquad
 \sum_{i=0}^{n-2}\binom{n-i}{2}2^i
    =2^{n+1}-\frac{n^2+3n+4}{2}.
 \tag{P-F28j}
\]
They follow by induction on $n$: the left sides
are twice their values at $n-1$, plus respectively
$n-1$ and $\binom n2$, and both identities begin
at $n=2$. The second sum in (P-F28i) is twice the
first sum of (P-F28j) with $n$ replaced by $n-1$.

Subtracting (P-F28h) from (P-F28i) yields
\[
 V_{\mathrm{cub}}-V_{\mathrm{ad}}
     =3M-\frac{n^2+3n+2}{2}
     =3M-\binom{n+2}{2}.
 \tag{P-F28k}
\]
This is (P-F28b). Each conormal eigenvalue is
nontrivial by P-F24, so its determinant factor
is nonzero. Combining it with $e_Y=1$ proves
(P-F28c).
\end{proof}

\paragraph{Role in the subsequent global argument.}
The local calculation does not compare possible boundary
components or establish their absence. P-F30 supplies
that separate support classification, and P-F31 then
verifies the global hypotheses of P-F23. Only at that
stage does the local coefficient computed here become
the actual normalized leading coefficient of the whole
Hilbert series. Thus the proof of the global pole uses
classification of the fixed locus, without an assumed
inequality between cyclotomic valuations of different
components.
